\documentclass[11pt]{article}
\usepackage{amsmath,amssymb,amsthm}
\usepackage[margin=2.5cm]{geometry}

\newtheorem{theorem}{Theorem}
\newtheorem{lemma}[theorem]{Lemma}
\newtheorem{corollary}[theorem]{Corollary}
\theoremstyle{remark}
\newtheorem{remark}[theorem]{Remark}
\theoremstyle{definition}
\newtheorem*{statement}{Official statement}

\title{Problem 4, Cell 2: Sun's conjecture for $k=4$}
\author{}
\date{}

\begin{document}
\maketitle

\begin{statement}
For integers $a$ and $m\ge1$ write $a\ (\mathrm{mod}\ m)=\{a+mt:t\in\mathbb{Z}\}$. A finite family
of classes $a_1\ (\mathrm{mod}\ m_1),\dots,a_k\ (\mathrm{mod}\ m_k)$ is pairwise disjoint if
$a_i\ (\mathrm{mod}\ m_i)\cap a_j\ (\mathrm{mod}\ m_j)=\emptyset$ whenever $i<j$. Nothing is assumed
about the moduli beyond $m_i\ge1$: they may repeat, and the residues $a_i$ may repeat as well. By the
Chinese remainder theorem,
\[
  a_i\ (\mathrm{mod}\ m_i)\cap a_j\ (\mathrm{mod}\ m_j)\neq\emptyset\iff\gcd(m_i,m_j)\mid a_i-a_j.\tag{$*$}
\]
\emph{Question.} Let $k\ge2$ and let $a_1\ (\mathrm{mod}\ m_1),\dots,a_k\ (\mathrm{mod}\ m_k)$ be
pairwise disjoint. Must there be a pair $i<j$ with $\gcd(m_i,m_j)\ge k$?

The same question for four classes. Prove the statement for $k=4$.
\end{statement}

\noindent Theorem~\ref{thm:main} below is exactly this statement. The notation $a\bmod m$ below is the class
$a\ (\mathrm{mod}\ m)$ of the statement. Lemma~\ref{lem:crt} is $(*)$, proved for completeness.

\paragraph{Known result.} The question is Sun's disjoint congruence classes conjecture~\cite{Sun06}.
O'Bryant~\cite{OBr07} proved it for all $k\le20$, so this case is known. The proof below is an
independent proof and is self-contained.

For an integer $m\ge1$ and $a\in\mathbb{Z}$ write $a\bmod m=\{x\in\mathbb{Z}: x\equiv a\pmod m\}$.
A family of residue classes is \emph{disjoint} if its members are pairwise disjoint.

\begin{theorem}\label{thm:main}
Let $a_1\bmod m_1,\dots,a_4\bmod m_4$ be four pairwise disjoint residue classes. Then
$\gcd(m_i,m_j)\ge4$ for some $i\ne j$.
\end{theorem}

\section{When do two classes meet?}

\begin{lemma}\label{lem:crt}
Let $m,n\ge1$, $a,b\in\mathbb{Z}$ and $g=\gcd(m,n)$. Then
$(a\bmod m)\cap(b\bmod n)\ne\emptyset$ if and only if $g\mid a-b$.
\end{lemma}

\begin{proof}
($\Rightarrow$) If $x\equiv a\pmod m$ and $x\equiv b\pmod n$, then $g\mid m$ and $g\mid n$ give
$a\equiv x\equiv b\pmod g$.

($\Leftarrow$) By B\'ezout there are $u,v\in\mathbb{Z}$ with $um+vn=g$. Write $a-b=gt$ and set
$x=a-umt$. Then $x\equiv a\pmod m$, and $x=a-(g-vn)t=b+vnt\equiv b\pmod n$.
\end{proof}

\begin{corollary}\label{cor:small}
If $a\bmod m$ and $b\bmod n$ are disjoint and $g=\gcd(m,n)$, then $g\ge2$ and $a\not\equiv b\pmod g$.
\end{corollary}

\begin{proof}
By Lemma~\ref{lem:crt}, disjointness means $g\nmid a-b$, and this forces $g\ne1$.
\end{proof}

\section{Proof of Theorem~\ref{thm:main}}

Suppose, for contradiction, that $\gcd(m_i,m_j)\le3$ for all $i\ne j$. Write $g_{ij}=\gcd(m_i,m_j)$
and $[4]=\{1,2,3,4\}$. By Corollary~\ref{cor:small},
\begin{equation}\label{eq:23}
  g_{ij}\in\{2,3\}\qquad\text{for all } i\ne j .
\end{equation}
Let
\[
  E=\{i\in[4]: 2\mid m_i\},\qquad T=\{i\in[4]: 3\mid m_i\}.
\]

\paragraph{Step 1: $E\cup T=[4]$ and $|E\cap T|\le1$.}
Given $i$, pick any $j\ne i$. By \eqref{eq:23}, $g_{ij}$ is $2$ or $3$ and divides $m_i$, so
$i\in E\cup T$. If two distinct indices $i,j$ were both in $E\cap T$, then $6\mid m_i$ and $6\mid m_j$.
This would give $g_{ij}\ge6$, contradicting \eqref{eq:23}.

\paragraph{Step 2: every pair lies in $E$ or in $T$.}
If $g_{ij}=2$, then $2\mid m_i$ and $2\mid m_j$, so $i,j\in E$. If $g_{ij}=3$, then $i,j\in T$.

\paragraph{Step 3: $E\setminus T=\emptyset$ or $T\setminus E=\emptyset$.}
Suppose $i\in E\setminus T$ and $j\in T\setminus E$. By Step~2, either $i,j\in E$, contradicting
$j\notin E$, or $i,j\in T$, contradicting $i\notin T$.

\paragraph{Step 4: conclusion.}
By Step~3 and Step~1, one of two cases holds.

\emph{Case A: $T\setminus E=\emptyset$.} Then $[4]=E\cup T=E$. If some pair had $g_{ij}=3$, Step~2
would put $i,j\in T\cap E$. That contradicts $|E\cap T|\le1$ from Step~1. So by \eqref{eq:23},
$g_{ij}=2$ for every pair $i\ne j$. Corollary~\ref{cor:small} then gives $a_i\not\equiv a_j\pmod2$
for all $i\ne j$. So the four integers $a_1,\dots,a_4$ would be pairwise incongruent modulo $2$,
which the pigeonhole principle forbids (only $2<4$ classes).

\emph{Case B: $E\setminus T=\emptyset$.} Then $[4]=T$. If some pair had $g_{ij}=2$, Step~2 would put
$i,j\in E\cap T$, again contradicting Step~1. So $g_{ij}=3$ for every pair. Corollary~\ref{cor:small}
then gives that $a_1,\dots,a_4$ are pairwise incongruent modulo $3$. This is impossible by
pigeonhole (only $3<4$ classes).

Both cases are contradictory. Hence $g_{ij}\ge4$ for some $i\ne j$. \qed

\begin{remark}
\begin{enumerate}
\item The bound is sharp: $r\bmod4$ for $r=0,1,2,3$ are disjoint and all pairwise gcds equal $4$.
\item The argument does not use the density bound $\sum_i 1/m_i\le1$. Only the arithmetic of
  the primes $2$ and $3$ is needed. Moduli $m_i=1$ are excluded automatically by \eqref{eq:23}.
\item \emph{Computer sanity check, not part of the proof.} The script \texttt{check\_p4.py} searched
  all moduli $m_1\le\dots\le m_4\le60$ with all pairwise gcds $\le3$ (214\,078 tuples), and all
  residues (with $a_1=0$, by translation). It found no disjoint system, which agrees with the theorem.
  Disjointness is checked literally, from the definition, and not via $(*)$: the classes $a\bmod m$
  and $b\bmod n$ meet iff $b-a\equiv mt\pmod n$ for some $t$, and the set $\{mt\bmod n:0\le t<n\}$ is
  listed explicitly. For moduli $\le20$ the search was repeated with an element-by-element
  comparison over one period.
\end{enumerate}
\end{remark}

\begin{thebibliography}{9}
\bibitem{OBr07} K.~O'Bryant, \emph{On Z.-W.\ Sun's disjoint congruence classes conjecture}, in
  \emph{Combinatorial Number Theory}, de Gruyter, Berlin, 2007, pp.~403--411; arXiv:math/0604347.
\bibitem{Sun06} Z.-W.~Sun, \emph{Finite covers of groups by cosets or subgroups}, Internat.\ J.\ Math.\
  \textbf{17} (2006), 1047--1064; arXiv:math/0501451.
\end{thebibliography}

\end{document}
